% Part: first-order-logic
% Chapter: introduction
% Section: syntax

\documentclass[../../../include/open-logic-section]{subfiles}

\begin{document}

\olfileid[pt]{fol}{int}{syn}

\olsection{Sintaxe}

Primeiro devemos tornar preciso quais cadeias de símbolos contam como
!!{sentence}s da lógica de primeira ordem.  Faremos isso mais adiante; por
enquanto, procederemos por exemplos.  Os blocos de construção básicos---o
vocabulário---da lógica de primeira ordem dividem-se em duas partes. A
primeira parte são os símbolos que usamos para dizer coisas específicas ou
para destacar coisas específicas. Destacamos coisas usando
!!{constant}s, e dizemos coisas sobre o que destacamos usando
!!{predicate}s.  Por exemplo, poderíamos usar $\Obj a$ como
!!a{constant} para destacar uma única coisa, e então dizer algo sobre ela
usando a !!{sentence}~$\Atom{\Obj P}{\Obj a}$.  Se você tem significados
para ``$\Obj a$'' e ``$\Obj P$'' em mente, pode ler $\Atom{\Obj P}{\Obj a}$
como uma sentença do inglês (e provavelmente já fez isso quando aprendeu
lógica formal pela primeira vez).  Uma vez que se tem tais !!{sentence}s
simples da lógica de primeira ordem, pode-se construir outras mais complexas
usando a segunda parte do vocabulário: os símbolos lógicos (conectivos e
quantificadores). Assim, por exemplo, podemos formar expressões como
$(\Atom{\Obj P}{\Obj a} \land \Atom{\Obj Q}{\Obj b})$ ou
~$\lexists[\Obj x][\Atom{\Obj P}{\Obj x}]$.

Para fornecer as definições precisas de semântica e as regras dos nossos
sistemas de !!{derivation} necessárias para o estudo meta-lógico rigoroso,
temos primeiro que dar uma definição precisa do que conta como
!!a{sentence} da lógica de primeira ordem. A ideia básica é fácil o
suficiente de entender: há algumas !!{sentence}s simples que podemos formar
apenas a partir de !!{predicate}s e !!{constant}s, como~$\Atom{\Obj P}{\Obj
a}$. E então a partir destas formamos outras mais complexas usando os
conectivos e quantificadores. Mas quais são exatamente as regras pelas quais
nos é permitido formar !!{sentence}s mais complexas?  Elas devem ser
especificadas; caso contrário, não definimos ``!!{sentence} da lógica de
primeira ordem'' com precisão suficiente. Há algumas questões. A primeira
é obter as cadeias certas para contarem como !!{sentence}s. A segunda é
fazê-lo de tal forma que possamos dar provas matemáticas sobre
\emph{todas} as !!{sentence}s. Finalmente, teremos também que dar definições
precisas de algumas operações rudimentares com !!{sentence}s, como
``substituir todo $\Obj x$ em~$!A$ por~$\Obj b$.'' O problema é que os
quantificadores e !!{variable}s que temos na lógica de primeira ordem
tornam não inteiramente óbvio como isso deve ser feito. Por exemplo,
$\lexists[\Obj x][\Atom{\Obj P}{\Obj a}]$ deve contar como
!!a{sentence}? E quanto a $\lexists[\Obj x][\lexists[\Obj
x][\Atom{\Obj P}{\Obj x}]]$? Qual deveria ser o resultado de ``substituir
$\Obj x$ por~$\Obj b$ em $(\Atom{\Obj P}{\Obj x} \land \lexists[\Obj
x][\Atom{\Obj P}{\Obj x}])$''?

\end{document}
